% Section 1 — Introduction.
% This file demonstrates the building blocks you will use most: blocks, lists, emphasis
% and citations. Delete the examples and write your own.

\begin{frame}{Motivation}

% A block with a heading. Use it to group one idea per block.
\begin{block}{The problem}
\begin{itemize}
    \item State the problem in one line that a non-specialist understands
    \item Say who is affected, and how much
    \item Resist the urge to explain your method here
\end{itemize}
\end{block}

% A block with no heading, written as {}, is a quiet container for supporting points.
\begin{block}{}
\begin{itemize}
    \item Supporting detail, with a citation~\cite{W3021074122}
    \item A second point, \textbf{emphasised} where it matters
\end{itemize}
\end{block}

\end{frame}


% [allowframebreaks] lets a long frame continue onto a second slide automatically.
% Beamer appends "I", "II" to the title. Useful while drafting; tighten before presenting.
\begin{frame}[allowframebreaks]{Background}

\begin{block}{What is already known}
\begin{itemize}
    \item One finding per line, each attached to its source~\cite{W3135286407}
    \item Keep lines short enough to read from the back of the room
\end{itemize}
\end{block}

\begin{block}{What is not known}
\begin{itemize}
    \item The gap your work fills, stated as plainly as you can manage
\end{itemize}
\end{block}

\end{frame}
